\section{Neural Compute, Memory y Personalization}

A continuación, el ponente compara al agent con un sistema de cómputo: el modelo funciona como una neural compute unit invocable, las invocaciones repetidas forman un ciclo, el scratchpad guarda el estado de trabajo y la memoria a largo plazo asume la persistencia entre sesiones. La analogía tiene mucho valor didáctico, pero hay que dejar claro que el Transformer no es una CPU ni el embedding store es un disco. La analogía sirve para separar el cómputo, la memoria de trabajo y el estado persistente.

\subsection{El modelo como Neural Compute Unit}

La primera figura asigna la secuencia máxima de tokens de entrada a la secuencia máxima de tokens de salida y subraya que una invocación del modelo es un cómputo dentro de una ventana finita. La segunda, la analogía con las instrucciones MIPS, indica lo siguiente: un procesador tradicional tiene una semántica de instrucciones fija, mientras que la «instrucción» de un modelo de lenguaje se interpreta a partir del lenguaje natural y del contexto en conjunto; por eso es flexible, pero inestable.

\lecturefigure{slide-46-neural-compute-unit.jpg}{Una invocación del modelo abstraída como neural compute unit}{00:28:00}

\lecturefigure{slide-47-mips-instruction-analogy.jpg}{Analogía entre las instrucciones fijas de MIPS y la invocación de un modelo de lenguaje natural}{00:28:34}

\begin{knowledgebox}{Correspondencias de la analogía}
Los tokens de entrada equivalen al programa y los datos; los parámetros del modelo, a una estructura de cómputo fija; los tokens de salida, al resultado del cómputo, y el scratchpad, al área de trabajo. Sin embargo, el modelo de lenguaje carece de la semántica estricta, las transacciones y el determinismo de una ISA de hardware; por ello, cada invocación necesita schema, validación y manejo de errores que lo compensen.
\end{knowledgebox}

\subsection{Looped Transformer: convertir una generación en un cómputo iterativo}

La figura del Looped Transformer devuelve la salida a la entrada y forma un ciclo entre scratchpad, memory e instructions. Los sistemas de agents hacen precisamente esto: mediante orchestration, extienden una sola generación hacia adelante a un algoritmo de varios pasos. Si el estado interno en la iteración $k$ es $z_k$, puede abstraerse como
\[
z_{k+1}=F_\theta(z_k,o_k,m_k),\qquad
a_k=G(z_{k+1}),
\]
donde $F_\theta$ es el cómputo del modelo, $o_k$ es la nueva observación, $m_k$ es la memoria recuperada y $G$ asigna el estado interno a una acción de herramienta. El sistema debe definir cuándo detenerse; de lo contrario, el propio ciclo se convierte en un modo de falla.

\lecturefigure{slide-48-looped-transformer.jpg}{Looped Transformer: scratchpad, memory e instructions entran en el circuito iterativo}{00:28:51}

\begin{importantbox}{La condición de paro importa más que «pensarlo un poco más»}
Las condiciones de terminación incluyen, como mínimo, la tarea completada, el presupuesto agotado, el estado repetido, el umbral de riesgo, la cancelación por parte del usuario y la intervención humana. Un reasoning loop sin condición de paro convierte la incertidumbre en tokens, latencia y costos de herramientas.
\end{importantbox}

\subsection{Long-term Memory: persistencia, recuperación, actualización y olvido}

La slide de memoria a largo plazo describe la memory como long-lived y persistent, y enumera embedding, retrieval models, jerarquía, consistencia temporal, estructura y adaptación en línea. La primera vez que aparecen hay que distinguir: la \term{working memory} es el estado de la tarea actual; la \term{long-term memory} es la información persistente entre sesiones; la \term{model memory} puede referirse al conocimiento aprendido en los parámetros. Las tres tienen mecanismos de actualización distintos.

\lecturefigure{slide-49-long-term-memory.jpg}{Mecanismos y problemas abiertos de la memoria a largo plazo}{00:35:17}

Una expresión básica de recuperación es
\[
m_t=\operatorname{TopK}_{i}\;\operatorname{sim}(q_t,e_i),
\]
donde $q_t$ es el embedding de la consulta actual, $e_i$ es la representación de la entrada de memoria y $m_t$ son las top-$k$ entradas devueltas. Recuperar solo por similitud trae contenido caduco o sensible; por eso, la puntuación real debe incorporar tiempo, fuente, permisos y confianza:
\[
\text{score}_i=\alpha\,\text{sim}(q_t,e_i)+\beta\,\text{recency}_i+
\eta\,\text{trust}_i-\lambda\,\text{privacyRisk}_i.
\]

\begin{warningbox}{Contaminación de la memoria y persistencia de errores}
Si un malentendido de una conversación se escribe en la memoria a largo plazo, se reforzará una y otra vez en el futuro. Antes de escribir, hay que distinguir entre lo que el usuario declaró explícitamente, lo que el sistema infirió y el estado temporal de la tarea; las preferencias de alto impacto deben poder ser consultadas, corregidas y eliminadas por el usuario.
\end{warningbox}

\begin{lstlisting}[caption={Registro de memoria con fuente y política de caducidad}]
memory_item = {
  "fact": normalized_content,
  "source": source_event,
  "confidence": confidence,
  "scope": ["travel", "restaurants"],
  "created_at": timestamp,
  "expires_at": expiry,
  "sensitivity": sensitivity,
  "user_editable": True
}
\end{lstlisting}

\subsection{La Personalization es un problema de alineación usuario--Agent}

La slide de personalización enumera a la vez preferencias explícitas e implícitas. Las explícitas, como alergias, asientos y platillos favoritos, suelen poder confirmarse directamente; las implícitas provienen de estadísticas de comportamiento y es más fácil que confundan una elección fortuita, una restricción de precio o un contexto histórico con un gusto estable. El objetivo del sistema puede escribirse como una optimización con restricciones:
\[
\max_\pi\;\mathbb{E}[U_{\text{user}}(\tau)]
\quad\text{s.t.}\quad
C_{\text{privacy}}(\tau)\le \epsilon_p,
\;C_{\text{safety}}(\tau)\le \epsilon_s.
\]
$U_{\text{user}}$ es la utilidad para el usuario, $C_{\text{privacy}}$ y $C_{\text{safety}}$ son los costos de privacidad y seguridad, y $\epsilon_p,\epsilon_s$ son los umbrales permitidos. La personalización no puede lograrse a costa de recopilar datos más allá de los límites.

\lecturefigure{slide-50-personalization.jpg}{Personalización: las preferencias explícitas e implícitas constituyen juntas la user-agent alignment}{00:36:42}

\lecturefigure{slide-51-personalization-challenges.jpg}{Retos de la personalización: recopilación de datos, aprendizaje, adaptación en línea y privacidad}{00:37:56}

Las dos figuras de personalización deben leerse junto con la de memoria a largo plazo: la memory decide «qué guardar», la personalization decide «cómo usarlo» y la privacy decide «cuándo no puede usarse». Por ejemplo, si un usuario eligió un asiento de pasillo para un viaje de negocios, eso no debe generalizarse de forma permanente como preferencia global; el sistema debe guardar la fuente, el contexto y la confianza, y volver a confirmar antes de una decisión de alto impacto. La personalización auténtica no consiste en que el agent adivine más, sino en que sepa qué preferencias son confiables, cuáles caducaron y cuáles deben preguntarse de nuevo.

\teachervoice{00:35:17--00:38:00, el ponente compara la memoria a largo plazo con un disco, define la personalización como user-agent alignment y enfatiza la consulta activa, el aprendizaje pasivo, SFT, human feedback, on-the-fly adaptation y privacy. Nota de clase: las preferencias aprendidas deben llevar fuente y un mecanismo para revocarlas.}

\begin{knowledgebox}{Términos clave: Memory y Personalization}
\begin{tabularx}{\textwidth}{>{\bfseries}p{0.2\textwidth} X X}
Término & Mecanismo central & Falla principal \\
\hline
Embedding memory & Recuperación de textos o eventos por similitud & Semánticamente cercano pero factualmente erróneo; caducidad temporal \\
Episodic memory & Guarda interacciones concretas y sus resultados & Mucho ruido; sensible en privacidad \\
Semantic memory & Condensa muchos eventos en hechos estables & Los errores de generalización quedan fijados a largo plazo \\
Online adaptation & Cambia la conducta al instante según la nueva retroalimentación & Deriva catastrófica; manipulación mediante retroalimentación maliciosa \\
Preference model & Predice las elecciones o la satisfacción del usuario & Confunde elecciones históricas con valores reales \\
\end{tabularx}
\end{knowledgebox}

\subsection{Resumen de la sección}

La analogía neural-compute ayuda a separar en capas la invocación del modelo, el ciclo, el estado de trabajo y la memoria a largo plazo, pero no sustituye el diseño real de interfaces. La memoria a largo plazo requiere recuperación, fuente, caducidad, permisos y edición por parte del usuario; la personalización requiere aprender bajo restricciones de utilidad, privacidad y seguridad, no recopilar datos de comportamiento sin límite.
